% ==============================================================================
% =================================== ABSTRACT =================================
% ==============================================================================

\chapter*{Abstract}
\addcontentsline{toc}{chapter}{Abstract}

Die Arbeit analysiert ein eskaliertes Klärungsgespräch zwischen einem \acs{RPA}/\acs{KI}-Entwickler und einer Kundin aus der Finanzabteilung. Anhand eines konstruierten Dialogskripts werden die fünf Watzlawick-Axiome, das Vier-Seiten-Modell nach Schulz von Thun sowie das neunstufige Eskalationsmodell nach Glasl angewendet.

Die Analyse zeigt, dass die Eskalation nicht auf der Sachebene (fehlerhaftes Modul vs. fehlende Zuarbeiten), sondern auf der Beziehungsebene entsteht: Die Kundin interpretiert den sachlichen Hinweis auf ihre Versäumnisse bei den Datenlieferungen als Angriff auf ihre fachliche Kompetenz. Die Eskalation folgt dabei exakt den drei Antriebsmechanismen nach Glasl – der wachsenden Streitpunktlawine, der Arena-Ausweitung und der pessimistischen Antizipation – sowie dem übergreifenden Prinzip der Selbstansteckung. Der Konflikt erreicht Stufe 5 (Gesichtsverlust) bis Stufe 6 (Drohstrategien).

Für die Praxis ergeben sich drei Erfolgsfaktoren: frühzeitiges Schnittstellen- und Erwartungsmanagement, Priorisierung der Beziehungsebene im Störungsfall sowie deeskalative Gesprächsführung mit aktiven Zuhörtechniken.

% ==============================================================================
% =================================== ABSTRACT =================================
% ==============================================================================